% Source: https://artofproblemsolving.com/community/c3h1969880_school_competition
% Stable identity: aops-topic-1969880-q03
Join the centre of the hexagon to its six vertices. This divides the hexagon into $6$ equilateral triangles of side $20$, each of area $\dfrac{\sqrt3}{4}\cdot20^2=100\sqrt3$. Assign each student to one triangle containing him (a student on a shared edge is assigned to just one of the triangles).

\begin{asy}
usepackage("amsmath");
size(4.5cm);
defaultpen(fontsize(10pt));
real r=2;
pair[] H; for(int i=0;i<6;++i) H.push(r*dir(60*i));
fill(H[0]--H[1]--(0,0)--cycle,palegray);
draw(H[0]--H[1]--H[2]--H[3]--H[4]--H[5]--cycle,linewidth(1));
for(int i=0;i<6;++i) draw((0,0)--H[i]);
label("$20$",(H[3]+H[4])/2,SW);
label("$100\sqrt{3}$",(H[0]+H[1])/3,fontsize(8pt));
label("\mbox{}",(min(currentpicture,user=true)+max(currentpicture,user=true))/2,black); // keeps the text after the figure black when LaTeX typesets the labels inline
\end{asy}

Since $19>6\cdot3$, by the pigeonhole principle some triangle $T$ contains at least $4$ students. Take four of them, $P,Q,R,S$. No three are collinear, so they form a (simple) quadrilateral: if their convex hull has $4$ vertices, take the convex quadrilateral; otherwise one point lies inside the triangle formed by the other three, and joining them in a suitable order gives a non-convex quadrilateral.

Its area is at most the area of the convex hull of $P,Q,R,S$, which is contained in $T$, so the area is at most $100\sqrt3$. Equality would require the convex hull to be all of $T$, so three of the points would be the vertices of $T$; the fourth point is then either on a side of $T$ (three collinear points, which is excluded) or strictly inside $T$, in which case the quadrilateral is non-convex and its area is the area of $T$ minus a triangle of positive area. Either way equality is impossible, so the quadrilateral has area less than $100\sqrt3\ \text{m}^2$. $\blacksquare$
